\section{条件数衡量问题本身的敏感性}
\label{sec:08-Numerical-Analysis/01-Floating-Point-and-Error/03}

你花了三天优化算法，逐行审查浮点运算顺序，把每次舍入都安排得无可挑剔。跑出来结果还是错的。

不是你的代码有问题。是问题本身在放大误差。就像你拿着世界上最好的听诊器去听一枚硬币掉在一百米外的声音——听诊器完美无瑕，但信号路上有一百米的衰减。

条件数衡量的是这条路的衰减（或放大）系数。它只跟问题定义有关，跟你怎么算无关。

\subsection{算法还没写，答案就可能已经不稳了}

把同一个计算要求交给两台机器、两种编程语言、两个独立的程序员。如果问题本身病态，三个结果很可能互不吻合——不是因为谁的程序有 bug，而是因为输入数据里各自进了不同的末位扰动，被问题放大成了截然不同的输出。

这最容易产生一种徒劳的调试循环：发现结果不准，怀疑代码，重构算法，结果差不多，再换个库，还是差不多。真正需要的是退一步问：在这个点上，输入稍微变一丁点，精确答案会变多少？如果变动超过你能接受的范围，那么无论算法多么稳定、精度多高，都无法从输入里挤出不存在的信息。

条件数描述的就是这个「输入扰动到输出扰动的放大倍数」——在给定输入点附近的一个局部性质。它与算法选择无关。

\subsection{绝对条件数与相对条件数}

从一阶 Taylor 展开出发。对于可微标量函数 \(y = f(x)\)，输入有微小扰动 \(\Delta x\) 时：
\[
|\Delta y| \approx |f'(x)| \cdot |\Delta x|.
\]
绝对条件数直接出现在这个关系里：
\[
\kappa_{\mathrm{abs}}(x) = |f'(x)|.
\]
它把输入的绝对扰动乘以导数的大小，得到输出的绝对扰动。

更常见的是相对条件数——输入变了百分之几，答案变百分之几：
\[
\frac{|\Delta y|}{|y|}
\approx
\left| \frac{x f'(x)}{f(x)} \right|
\cdot
\frac{|\Delta x|}{|x|}.
\]
所以
\[
\kappa_{\mathrm{rel}}(x)
= \left| \frac{x f'(x)}{f(x)} \right|.
\]

选择绝对还是相对，不是形式上的偏好，而是你在意什么。测量一栋楼的高度，在意的是厘米级的绝对误差；计算两个近等大数的微小差值，在意的是相对误差——后者更容易因条件数爆炸而失控。

\subsubsection{一个接近根时条件数爆炸的简单例子}

取 \(f(x) = x - 1\)。它的绝对条件数恒为 1——输入变多少，输出就变多少，体面得很。

但相对条件数是
\[
\kappa_{\mathrm{rel}}(x) = \left| \frac{x}{x - 1} \right|.
\]
当 \(x\) 逼近 1 时，分母趋于零，条件数飞向无穷。原因很直白：我们试图从两个极其接近的数（\(x\) 和 \(1\)）之差读出结果。输入的微小相对扰动变成差值的巨大相对扰动。

上一节讨论的灾难性消去，从条件数的角度看就有了更清晰的分类：
\begin{itemize}
\tightlist
\item 如果两个相近数来自独立测量，差值的计算本身就是病态问题——改写表达式无能为力。
\item 如果两个相近数来自同一个解析表达式，代数改写可以利用内部结构绕开这个病态步骤。比如把 \(1 - \cos x\) 写成 \(2\sin^2(x/2)\)，是把一个条件数巨大的表达式拆解成了直接计算小量，不再经过相近数相减。
\end{itemize}

条件数帮你判断：这个损失是问题「本来就该有的」，还是你选的表达式「没必要的」。

\subsection{线性系统的条件数}

线性方程组 \(Ax = b\) 是更丰富的舞台。

先只扰动右端 \(b\)：由 \(A \Delta x = \Delta b\) 得 \(\Delta x = A^{-1} \Delta b\)。取与向量范数相容的矩阵算子范数（满足 \(\|Av\| \le \|A\| \|v\|\)），有
\[
\|\Delta x\| \le \|A^{-1}\| \|\Delta b\|.
\]
又从 \(\|b\| \le \|A\| \|x\|\) 得到相对误差的界：
\[
\frac{\|\Delta x\|}{\|x\|}
\le
\|A\| \|A^{-1}\|
\frac{\|\Delta b\|}{\|b\|}.
\]
中间那个乘积就是矩阵的条件数：
\[
\kappa(A) = \|A\| \|A^{-1}\|.
\]

在 2-范数下，它可以写得非常直观：
\[
\kappa_2(A) = \frac{\sigma_{\max}(A)}{\sigma_{\min}(A)}.
\]
这里 \(\sigma_{\max}\) 和 \(\sigma_{\min}\) 分别是 \(A\) 的最大和最小奇异值。

几何上，\(A\) 把单位球映成一个超椭球。最长半轴是 \(\sigma_{\max}\)，最短半轴是 \(\sigma_{\min}\)。条件数是两轴之比。如果最短半轴接近零——存在某个输入方向几乎被 \(A\) 压扁——那么求逆时这个方向会被放大 \(1/\sigma_{\min}\) 倍。矩阵越接近奇异，这个数字越大。

正交矩阵的 2-范数条件数是 1——它不改变任何方向的长度，只是旋转或反射，对误差既不放大也不压缩。奇异矩阵的条件数无定义（或者说无穷）——至少有一个方向的信息被彻底摧毁，求逆无从谈起。

\subsection{小残差不等于小误差}

你解 \(Ax = b\) 得到了一个近似解 \(\hat x\)。计算残差：
\[
r = b - A \hat x.
\]
如果 \(\|r\|\) 小到接近机器精度，你是不是已经快要摸到真解了？

不一定。残差小的含义是：\(\hat x\) 精确地解了一个右端略微改变的系统 \(A \hat x = b - r\)——这是一种后向误差（下一节会正式定义）。但真解与近似解之间的前向误差是
\[
x - \hat x = A^{-1} r.
\]
如果 \(A\) 病态，\(A^{-1}\) 会把一个微小的残差放大成巨大的解误差。看到 Solver 报告残差范数约 \(10^{-16}\)，只能说明算法实现很稳定；不能说明答案接近真解。要对前向误差有信心，你得把残差乘以条件数。

\subsection{最坏方向与实际方向}

条件数给出的是局部最坏方向的放大倍数。如果实际扰动恰好不跟最坏方向重合——这在实际数据中经常发生——实际误差可能远小于条件数给出的上界。这不是条件数「算错了」，而是它本来的定义就是一个上界。

这也引出一个重要的扩展。扰动有时带着结构：矩阵的对称性必须保持、非零模式不能改变、元素作为物理参数有天然的量级关系。此时允许每个元素各自独立随机变化的结构无关条件数过于保守。需要研究结构化条件数——只允许扰动落在符合问题约束的子空间内。

缩放变量和方程可以改变表示下的条件数数值，改善某些算法的行为。但真实的不可辨识性无法被缩放消除：如果两个模型参数对观测产生几乎完全相同的影响，换一套单位不会凭空造出辨识信息来。

\subsection{条件数告诉你能信几位}

把前面几节串起来，得到一个实用口诀。

假设输入数据的相对误差大致在 \(u\) 量级（比如 double precision 下 \(u \approx 10^{-16}\)），问题的条件数是 \(\kappa\)。那么即使算法每一步都是后向稳定的，最终输出的相对误差也常在 \(\kappa u\) 的量级上。

十进制下的翻译：条件数约 \(10^k\)，意味着你大概会丢掉 \(k\) 位有效数字。输入有 16 位，条件数 \(10^6\)，能信的大概只剩 10 位。

这当然只是量级判断。它需要 \(\kappa u \ll 1\) 以及局部线性近似合理。如果 \(\kappa u\) 已经到了 0.1 甚至更大，线性分析给出的是「你可能一位都不剩」——此时继续在算法细节上死磕不如回头增加输入精度、重新参数化，或者干脆换一个对输入不那么敏感的问题提法。